\newpage
\section{Results}\label{sec:results}
This section presents the empirical results from the benchmark described in Section \ref{sec:measurement}. The results cover handshake latency, message sizes, and memory footprint for each of the nine evaluated variants: classical WireGuard, four ML-KEM variants (hybrid and pure, at security levels 3 and 5), and four HQC variants (hybrid and pure, at security levels 3 and 5). The results are then used to answer the research questions posed in Section \ref{research_questions} and to draw comparisons with the related work.


\subsection{Metrics}
Table ref{tab:Results} reports the following values for each variant:

\begin{itemize} 
    \item \textbf{Initial [Bytes]:} The on-wire size of the initiation message in bytes. 
    \item \textbf{Response [Bytes]:} The on-wire size of the response message in bytes. 
    \item \textbf{Baseline RTT [ms]:} The mean RTT of five consecutive pings issued on an already-established session, representing pure network overhead with no handshake. The standard deviation of the baseline is shown in parentheses. \item \textbf{Handshake Time [ms]:} The mean baseline-subtracted RTT over 5000 trials, reflecting the duration of the handshake itself. The standard deviation is shown in parentheses. 
\end{itemize}

The baseline RTT is consistent at approximately \qty{0.5}{ms} across all nine variants, which confirms that the underlying virtual network conditions did not change between runs and that differences in handshake time reflect cryptographic overhead rather than measurement noise.

\begin{table}[H]
\centering
\begin{tabular}{|l|rrrr|}
\hline
\multirow{2}{*}{\textbf{Variant}} & Initial & Response & Baseline RTT & Handshake Time \\
                         & [Bytes] & [Bytes] & [ms] & [ms] \\
\hline
\multirow{2}{*}{Classic} & 148 & 92 & 0.522 & 1.820 \\
                         &  &  & (0.041) & (0.245) \\
\hline
\multirow{2}{*}{Pure-MLKEM-768} & 3540 & 1088 & 0.511 & 1.172 \\
                         &  &  & (0.034) & (0.231) \\
\hline
\multirow{2}{*}{Pure-MLKEM-1024} & 4788 & 1628 & 0.521 & 1.437 \\
                         &  &  & (0.052) & (0.243) \\
\hline
\multirow{2}{*}{Hybrid-MLKEM-768} & 1332 & 1180 & 0.525 & 2.228 \\
                         &  &  & (0.063) & (0.283) \\
\hline
\multirow{2}{*}{Hybrid-MLKEM-1024} & 1716 & 1660 & 0.517 & 2.374 \\
                         &  &  & (0.035) & (0.336) \\
\hline
\multirow{2}{*}{Pure-HQC-192} & 18106 & 9038 & 0.529 & 364.038 \\
                         &  &  & (0.047) & (14.302) \\
\hline
\multirow{2}{*}{Pure-HQC-256} & 28995 & 14481 & 0.520 & 668.182 \\
                         &  &  & (0.044) & (23.451) \\
\hline
\multirow{2}{*}{Hybrid-HQC-192} & 4670 & 9070 & 0.549 & 201.687 \\
                         &  &  & (0.059) & (8.220) \\
\hline
\multirow{2}{*}{Hybrid-HQC-256} & 7393 & 14513 & 0.538 & 365.557 \\
                         &  &  & (0.062) & (14.777) \\
\hline
\end{tabular}
\caption{Handshake benchmark results. Values in parentheses are standard deviations. Handshake time is the mean baseline-subtracted RTT over 5000 trials.}
\label{tab:Results}
\end{table}


\subsection{Message Size}
HQC very large ciphertexts


\subsection{Handshake Latency}
\subsubsection{ML-KEM}
ML-KEM faster than classic due to AES-NI on CPU

\subsubsection{HQC}
HQC very slow, probably due to vSwitch using software fragmentation reassembly on the single-core processor.


\subsection{Memory and Storage}



\subsection{Comparisons With Related Work}\label{sec:results-comparisons}

\subsubsection{PQ-WireGuard}


\subsubsection{Rosenpass}


\subsubsection{Rebar}


\subsection{Research Questions}